\section{Discussion and limitations}

Optimizer response can distinguish factorizations that have the same
effective weights.
For positive full-rank softmax/shared-basis factors, complete Euclidean
response identifies the current coordinates up to permutation on a common
product fiber. Explicit inverse bounds extend this observation to noisy
products and finitely many matrix-valued queries on declared nondegenerate
domains. A decision can require less information: balanced
Hilbert--Schmidt-response-optimal tying only needs off-diagonal router Gram
entries. Thus response-optimal tying can differ from weight clustering while
requiring less than full factor recovery.

The identification guarantees require known positive Euclidean rates and
temperature, internal access to the effective product and responses, and
rank/interior margins; the specified finite-query construction additionally
requires $D-K\ge L$ and a domain-dependent product-noise threshold.
Constants are conservative, query optimality is unproved, and no efficient
globally feasible recovery algorithm is supplied. Empirical rejection rules cannot certify these assumptions, and the
experiments include false acceptances. The results do not identify an AdamW trajectory, historical
sharing graph, or nonlinear output-mixture model.

Decision equivalence has narrower conditions than identification. The
response objective uses equal group capacities, matching hard/soft basis
rates, fixed hard routes and an isotropic full-operator error. It need not
minimize error for observed task gradients or an adaptive optimizer.
The trajectory bound requires response control throughout an interval,
which a single-checkpoint optimum does not provide. The proof and experiments establish that crossings are possible. Their
frequency and attraction properties under natural training remain open:
the evidence comes from a constructed loss, a selected real-model state,
and a finite exploratory grid. In particular, ordinary SGD produced no
crossings in the grid with its unretuned learning-rate schedule.

The recovery experiments cover a fixed-rank 64-channel additive module in
nine dense pretrained models up to 8B. The trajectory grid covers three
models up to 0.6B with 12 selected projections; local refinement uses one
selected SmolLM2 state and one fixed real batch. These are frozen-backbone
adapter studies, not full-model training or architectural universality.
Model-family comparisons confound architecture, pretraining and, for the
selected Qwen models, post-training. In the original matched decision test, factor recovery provided no
incremental benefit over effective-weight clustering. The later crossing
experiments address how decisions change during training; downstream
performance and compression gains remain unestablished. The results separate
the information needed to identify current factors from that needed for a
particular decision, and establish conditions under which training can move
the decisions apart.
